% ===========================================================================
\section{Verificación}
\label{sec:verificacion}
% ===========================================================================

Los conjuntos de Primeros se calcularon a mano y se contrastaron con una herramienta escrita para
este trabajo, \texttt{herramientas/primeros.py}. La herramienta existe para \textbf{contrastar} la
cuenta manual, no para reemplazarla: cuando las dos difirieron, se resolvió primero cuál de las dos
estaba equivocada.

\subsection{Condición LL(1)}

\begin{consola}
$ python3 herramientas/primeros.py gramatica/minijava-base.g
--- 60 no terminales, 126 producciones, 51 terminales ---
NO ES LL(1): 16 conflicto(s)

$ python3 herramientas/primeros.py gramatica/minijava-ll1-logros.g
--- 96 no terminales, 201 producciones, 60 terminales ---
ES LL(1): toda produccion de todo no terminal tiene Primeros disjuntos.
\end{consola}

\subsection{Atomicidad de los logros}

La condición de la sección~\ref{sec:atomicidad} exige que la gramática sea LL(1) para
\emph{todo} subconjunto de logros, no sólo con todos prendidos. Cada producción lleva su marca y la
herramienta recorre las $2^{10} = 1024$ configuraciones:

\begin{consola}
$ python3 herramientas/primeros.py gramatica/minijava-ll1-logros.g --combinatoria
10 logros -> 1024 configuraciones

LAS 1024 CONFIGURACIONES SON LL(1).
Cada logro se puede prender y apagar de forma independiente.
\end{consola}

\subsection{Ambigüedad frente a conflicto de predicción}

Para distinguir los dos defectos de la sección~\ref{sec:dos-defectos} se comparan los rangos de
longitud derivable: si no se solapan, los lenguajes son disjuntos y no puede haber ambigüedad.

\begin{consola}
$ python3 herramientas/primeros.py gramatica/minijava-base.g \
      --longitudes VarEncadenada MetodoEncadenado
<VarEncadenada>: entre 2 y 2 tokens
<MetodoEncadenado>: entre 4 y inf tokens

Los rangos NO se solapan: ninguna cadena puede estar en los dos
lenguajes, asi que <VarEncadenada> y <MetodoEncadenado> NO son ambiguos
entre si. Si comparten prefijo izquierdo es un conflicto de PREDICCION
(LL(1)), no de ambiguedad, y se arregla factorizando.
\end{consola}

\subsection{Casos de aceptación y rechazo}

La herramienta incluye un simulador del esquema de traducción --- elige la producción cuyo cuerpo
tiene al token actual en sus Primeros, y si ninguna lo tiene pero el no terminal es vaciable, no
hace nada --- que permite probar la gramática \textbf{antes} de escribir el analizador. Se
verifican \textbf{54 casos}: por cada logro, la sintaxis que debe aceptar y la que debe rechazar
cuando está apagado.

\begin{consola}
$ bash herramientas/casos-gramatica.sh
--- nucleo (sin ningun logro) ---
  ok    clase vacia
  ok    if/else anidado
  ok    encadenamiento
  ok    a = b = c NO es valido
--- MiniLambdas ---
  ok    x -> x            (1 param sin parentesis)
  ok    ( ) -> 1          (0 params)
  ok    ( a ) -> a        (1 param con parentesis)
  ok    ( a , b ) -> a    (2 params)
  ok    ( a ) + 1         (parentizada, NO lambda)
  ok    cuerpo con llaves NO se permite
  ok    parametro con tipo NO se permite
--- Genericidad Avanzada ---
  ok    List < List < Persona > > x ;  (anidado)
  ok    new List < > ( )               (diamante)
  ok    List < > x ;    diamante en posicion de TIPO
  ok    new List < > [ 3 ]  diamante con arreglo
...
TODOS LOS CASOS PASAN.
\end{consola}

Los casos que fijan las decisiones difíciles de este informe son, entre otros: \texttt{a = b = c}
se rechaza (no asociatividad de la asignación); \texttt{( a ) + 1} se acepta como expresión
parentizada y \texttt{( a ) -> a} como lambda, decidiéndose con un solo token; \texttt{List<> x;}
se rechaza y \texttt{new List<>()} se acepta.

\subsection{Validación contra los casos de la cátedra}

La verificación más fuerte no usa casos propios. La herramienta
\texttt{aceptar\_fuente.py} tokeniza un fuente \textbf{real} con el analizador léxico de la
Etapa 1 y le pasa la secuencia de tokens al simulador, de modo que la gramática se puede probar
contra los ocho casos entregados con los testers de esta etapa antes de que exista el analizador
sintáctico.

\begin{consola}
$ python3 herramientas/aceptar_fuente.py \
      ../vendor/testersSintactico/resources/*/*.java
  ok    sintCorrecto01.java      acepta
  ok    sintCorrecto02.java      acepta
  ok    sintCorrecto03.java      acepta
  ok    sintCorrecto04.java      acepta
  ok    sintError01.java         rechaza  [Error:,|3] ==
        se esperaba `{` y se encontro `,`
  ok    sintError02.java         rechaza  [Error:clasesita|3] ==
        se esperaba `idClase` y se encontro `idMetVar`
  ok    sintError03.java         rechaza  [Error:hola|5] ==
        se esperaba `;` y se encontro `idMetVar`
  ok    sintError04.java         rechaza  [Error:else|7] ==
        se esperaba `}` y se encontro `else`

TODOS LOS CASOS PASAN.
\end{consola}

Los cuatro casos válidos se aceptan y los cuatro inválidos se rechazan. Además ---y esto es lo
que importa para la corrección--- el rechazo se produce \textbf{en el token exacto} que declara la
expectativa de cada caso: el \texttt{==} indica que el código \texttt{[Error:lexema|línea]}
derivado de la gramática coincide con el que la cátedra espera en los cuatro. La gramática no sólo
reconoce el lenguaje correcto: falla donde tiene que fallar.

\subsection{Resumen}

\begin{center}
\begin{tabular}{@{}lrr@{}}
\toprule
 & \textbf{Gramática base} & \textbf{Gramática final} \\
\midrule
No terminales            & 60  & 96  \\
Producciones             & 126 & 201 \\
Terminales               & 51  & 60  \\
Ambigüedades             & 2   & 0   \\
Conflictos LL(1)& 16  & 0   \\
\bottomrule
\end{tabular}
\end{center}

El crecimiento se reparte entre la corrección del núcleo (seis no terminales) y las nueve
extensiones por logro. Ninguna ambigüedad y ningún conflicto de predicción quedan pendientes, en
ninguna de las 1024 configuraciones.

\subsection{Tabla de Primeros}

Conjunto de Primeros de cada producción de la gramática final, con todos los logros prendidos.

\begingroup\small
\setlength{\parindent}{0pt}
\setlength{\parskip}{0.7ex}
\raggedright
\hyphenpenalty=10000 \exhyphenpenalty=10000 \hbadness=10000
\par\hangindent=2.2em\hangafter=1
\regla{Inicial}{\nt{ListaClases} \tl{eof}}\nobreak\par
{\leftskip=1.6em\relax\noindent\textup{Primeros} $=$ \{ \tl{class}, \tl{eof}, \tl{interface} \}\par}
\par\hangindent=2.2em\hangafter=1
\regla{ListaClases}{\nt{Clase} \nt{ListaClases}}\nobreak\par
{\leftskip=1.6em\relax\noindent\textup{Primeros} $=$ \{ \tl{class} \}\par}
\par\hangindent=2.2em\hangafter=1
\regla{ListaClases}{\nt{Interfaz} \nt{ListaClases}}\nobreak\par
{\leftskip=1.6em\relax\noindent\textup{Primeros} $=$ \{ \tl{interface} \}\par}
\par\hangindent=2.2em\hangafter=1
\regla{ListaClases}{\eps}\nobreak\par
{\leftskip=1.6em\relax\noindent\textup{Primeros} $=$ \{ \eps \}\par}
\par\hangindent=2.2em\hangafter=1
\regla{Clase}{\tl{class} \tl{idClase} \nt{GenericidadOpcional} \nt{HerenciaOpcional} \tl{\{} \nt{ListaMiembros} \tl{\}}}\nobreak\par
{\leftskip=1.6em\relax\noindent\textup{Primeros} $=$ \{ \tl{class} \}\par}
\par\hangindent=2.2em\hangafter=1
\regla{Interfaz}{\tl{interface} \tl{idClase} \nt{GenericidadOpcional} \nt{ExtensionOpcional} \tl{\{} \nt{ListaMetodosInterfaz} \tl{\}}}\nobreak\par
{\leftskip=1.6em\relax\noindent\textup{Primeros} $=$ \{ \tl{interface} \}\par}
\par\hangindent=2.2em\hangafter=1
\regla{GenericidadOpcional}{\tl{<} \tl{idGen} \tl{>}}\nobreak\par
{\leftskip=1.6em\relax\noindent\textup{Primeros} $=$ \{ \tl{<} \}\par}
\par\hangindent=2.2em\hangafter=1
\regla{GenericidadOpcional}{\eps}\nobreak\par
{\leftskip=1.6em\relax\noindent\textup{Primeros} $=$ \{ \eps \}\par}
\par\hangindent=2.2em\hangafter=1
\regla{HerenciaOpcional}{\tl{extends} \nt{TipoReferencia}}\nobreak\par
{\leftskip=1.6em\relax\noindent\textup{Primeros} $=$ \{ \tl{extends} \}\par}
\par\hangindent=2.2em\hangafter=1
\regla{HerenciaOpcional}{\tl{implements} \nt{TipoReferencia}}\nobreak\par
{\leftskip=1.6em\relax\noindent\textup{Primeros} $=$ \{ \tl{implements} \}\par}
\par\hangindent=2.2em\hangafter=1
\regla{HerenciaOpcional}{\eps}\nobreak\par
{\leftskip=1.6em\relax\noindent\textup{Primeros} $=$ \{ \eps \}\par}
\par\hangindent=2.2em\hangafter=1
\regla{ExtensionOpcional}{\tl{extends} \nt{TipoReferencia}}\nobreak\par
{\leftskip=1.6em\relax\noindent\textup{Primeros} $=$ \{ \tl{extends} \}\par}
\par\hangindent=2.2em\hangafter=1
\regla{ExtensionOpcional}{\eps}\nobreak\par
{\leftskip=1.6em\relax\noindent\textup{Primeros} $=$ \{ \eps \}\par}
\par\hangindent=2.2em\hangafter=1
\regla{ListaMiembros}{\nt{Miembro} \nt{ListaMiembros}}\nobreak\par
{\leftskip=1.6em\relax\noindent\textup{Primeros} $=$ \{ \tl{boolean}, \tl{char}, \tl{float}, \tl{idClase}, \tl{idGen}, \tl{int}, \tl{private}, \tl{public}, \tl{static}, \tl{void} \}\par}
\par\hangindent=2.2em\hangafter=1
\regla{ListaMiembros}{\eps}\nobreak\par
{\leftskip=1.6em\relax\noindent\textup{Primeros} $=$ \{ \eps \}\par}
\par\hangindent=2.2em\hangafter=1
\regla{ListaMetodosInterfaz}{\nt{MetodoInterfaz} \nt{ListaMetodosInterfaz}}\nobreak\par
{\leftskip=1.6em\relax\noindent\textup{Primeros} $=$ \{ \tl{boolean}, \tl{char}, \tl{float}, \tl{idClase}, \tl{idGen}, \tl{int}, \tl{public}, \tl{void} \}\par}
\par\hangindent=2.2em\hangafter=1
\regla{ListaMetodosInterfaz}{\eps}\nobreak\par
{\leftskip=1.6em\relax\noindent\textup{Primeros} $=$ \{ \eps \}\par}
\par\hangindent=2.2em\hangafter=1
\regla{Miembro}{\nt{VisibilidadOpcional} \nt{MiembroSinVisibilidad}}\nobreak\par
{\leftskip=1.6em\relax\noindent\textup{Primeros} $=$ \{ \tl{boolean}, \tl{char}, \tl{float}, \tl{idClase}, \tl{idGen}, \tl{int}, \tl{private}, \tl{public}, \tl{static}, \tl{void} \}\par}
\par\hangindent=2.2em\hangafter=1
\regla{Miembro}{\nt{MiembroSinVisibilidad}}\nobreak\par
{\leftskip=1.6em\relax\noindent\textup{Primeros} $=$ \{ \tl{boolean}, \tl{char}, \tl{float}, \tl{idClase}, \tl{idGen}, \tl{int}, \tl{static}, \tl{void} \}\par}
\par\hangindent=2.2em\hangafter=1
\regla{Miembro}{\nt{Constructor}}\nobreak\par
{\leftskip=1.6em\relax\noindent\textup{Primeros} $=$ \{ \tl{public} \}\par}
\par\hangindent=2.2em\hangafter=1
\regla{VisibilidadOpcional}{\tl{public}}\nobreak\par
{\leftskip=1.6em\relax\noindent\textup{Primeros} $=$ \{ \tl{public} \}\par}
\par\hangindent=2.2em\hangafter=1
\regla{VisibilidadOpcional}{\tl{private}}\nobreak\par
{\leftskip=1.6em\relax\noindent\textup{Primeros} $=$ \{ \tl{private} \}\par}
\par\hangindent=2.2em\hangafter=1
\regla{VisibilidadOpcional}{\eps}\nobreak\par
{\leftskip=1.6em\relax\noindent\textup{Primeros} $=$ \{ \eps \}\par}
\par\hangindent=2.2em\hangafter=1
\regla{Constructor}{\tl{public} \tl{idClase} \nt{CuerpoDeMetodo}}\nobreak\par
{\leftskip=1.6em\relax\noindent\textup{Primeros} $=$ \{ \tl{public} \}\par}
\par\hangindent=2.2em\hangafter=1
\regla{MiembroSinVisibilidad}{\tl{static} \nt{TipoMetodo} \nt{Metodo}}\nobreak\par
{\leftskip=1.6em\relax\noindent\textup{Primeros} $=$ \{ \tl{static} \}\par}
\par\hangindent=2.2em\hangafter=1
\regla{MiembroSinVisibilidad}{\tl{void} \nt{Metodo}}\nobreak\par
{\leftskip=1.6em\relax\noindent\textup{Primeros} $=$ \{ \tl{void} \}\par}
\par\hangindent=2.2em\hangafter=1
\regla{MiembroSinVisibilidad}{\nt{TipoPrimitivo} \nt{DimensionesOpcionales} \tl{idMetVar} \nt{AtributoOMetodo}}\nobreak\par
{\leftskip=1.6em\relax\noindent\textup{Primeros} $=$ \{ \tl{boolean}, \tl{char}, \tl{float}, \tl{int} \}\par}
\par\hangindent=2.2em\hangafter=1
\regla{MiembroSinVisibilidad}{\tl{idGen} \nt{DimensionesOpcionales} \tl{idMetVar} \nt{AtributoOMetodo}}\nobreak\par
{\leftskip=1.6em\relax\noindent\textup{Primeros} $=$ \{ \tl{idGen} \}\par}
\par\hangindent=2.2em\hangafter=1
\regla{MiembroSinVisibilidad}{\tl{idClase} \nt{MiembroDesdeIdClase}}\nobreak\par
{\leftskip=1.6em\relax\noindent\textup{Primeros} $=$ \{ \tl{idClase} \}\par}
\par\hangindent=2.2em\hangafter=1
\regla{MiembroDesdeIdClase}{\nt{CuerpoDeMetodo}}\nobreak\par
{\leftskip=1.6em\relax\noindent\textup{Primeros} $=$ \{ \tl{(} \}\par}
\par\hangindent=2.2em\hangafter=1
\regla{MiembroDesdeIdClase}{\nt{TipoGenericoOpcional} \nt{DimensionesOpcionales} \tl{idMetVar} \nt{AtributoOMetodo}}\nobreak\par
{\leftskip=1.6em\relax\noindent\textup{Primeros} $=$ \{ \tl{<}, \tl{[}, \tl{idMetVar} \}\par}
\par\hangindent=2.2em\hangafter=1
\regla{AtributoOMetodo}{\nt{InicializacionDeAtributo} \tl{;}}\nobreak\par
{\leftskip=1.6em\relax\noindent\textup{Primeros} $=$ \{ \tl{;}, \tl{=} \}\par}
\par\hangindent=2.2em\hangafter=1
\regla{AtributoOMetodo}{\nt{CuerpoDeMetodo}}\nobreak\par
{\leftskip=1.6em\relax\noindent\textup{Primeros} $=$ \{ \tl{(} \}\par}
\par\hangindent=2.2em\hangafter=1
\regla{InicializacionDeAtributo}{\tl{=} \nt{ExpresionCondicional}}\nobreak\par
{\leftskip=1.6em\relax\noindent\textup{Primeros} $=$ \{ \tl{=} \}\par}
\par\hangindent=2.2em\hangafter=1
\regla{InicializacionDeAtributo}{\eps}\nobreak\par
{\leftskip=1.6em\relax\noindent\textup{Primeros} $=$ \{ \eps \}\par}
\par\hangindent=2.2em\hangafter=1
\regla{Metodo}{\tl{idMetVar} \nt{CuerpoDeMetodo}}\nobreak\par
{\leftskip=1.6em\relax\noindent\textup{Primeros} $=$ \{ \tl{idMetVar} \}\par}
\par\hangindent=2.2em\hangafter=1
\regla{CuerpoDeMetodo}{\nt{ArgsFormales} \nt{Bloque}}\nobreak\par
{\leftskip=1.6em\relax\noindent\textup{Primeros} $=$ \{ \tl{(} \}\par}
\par\hangindent=2.2em\hangafter=1
\regla{MetodoInterfaz}{\nt{PublicOpcional} \nt{TipoMetodo} \tl{idMetVar} \nt{ArgsFormales} \tl{;}}\nobreak\par
{\leftskip=1.6em\relax\noindent\textup{Primeros} $=$ \{ \tl{boolean}, \tl{char}, \tl{float}, \tl{idClase}, \tl{idGen}, \tl{int}, \tl{public}, \tl{void} \}\par}
\par\hangindent=2.2em\hangafter=1
\regla{PublicOpcional}{\tl{public}}\nobreak\par
{\leftskip=1.6em\relax\noindent\textup{Primeros} $=$ \{ \tl{public} \}\par}
\par\hangindent=2.2em\hangafter=1
\regla{PublicOpcional}{\eps}\nobreak\par
{\leftskip=1.6em\relax\noindent\textup{Primeros} $=$ \{ \eps \}\par}
\par\hangindent=2.2em\hangafter=1
\regla{TipoMetodo}{\nt{Tipo}}\nobreak\par
{\leftskip=1.6em\relax\noindent\textup{Primeros} $=$ \{ \tl{boolean}, \tl{char}, \tl{float}, \tl{idClase}, \tl{idGen}, \tl{int} \}\par}
\par\hangindent=2.2em\hangafter=1
\regla{TipoMetodo}{\tl{void}}\nobreak\par
{\leftskip=1.6em\relax\noindent\textup{Primeros} $=$ \{ \tl{void} \}\par}
\par\hangindent=2.2em\hangafter=1
\regla{Tipo}{\nt{TipoBase} \nt{DimensionesOpcionales}}\nobreak\par
{\leftskip=1.6em\relax\noindent\textup{Primeros} $=$ \{ \tl{boolean}, \tl{char}, \tl{float}, \tl{idClase}, \tl{idGen}, \tl{int} \}\par}
\par\hangindent=2.2em\hangafter=1
\regla{TipoBase}{\nt{TipoPrimitivo}}\nobreak\par
{\leftskip=1.6em\relax\noindent\textup{Primeros} $=$ \{ \tl{boolean}, \tl{char}, \tl{float}, \tl{int} \}\par}
\par\hangindent=2.2em\hangafter=1
\regla{TipoBase}{\nt{TipoReferencia}}\nobreak\par
{\leftskip=1.6em\relax\noindent\textup{Primeros} $=$ \{ \tl{idClase} \}\par}
\par\hangindent=2.2em\hangafter=1
\regla{TipoBase}{\tl{idGen}}\nobreak\par
{\leftskip=1.6em\relax\noindent\textup{Primeros} $=$ \{ \tl{idGen} \}\par}
\par\hangindent=2.2em\hangafter=1
\regla{DimensionesOpcionales}{\tl{[} \tl{]} \nt{DimensionesOpcionales}}\nobreak\par
{\leftskip=1.6em\relax\noindent\textup{Primeros} $=$ \{ \tl{[} \}\par}
\par\hangindent=2.2em\hangafter=1
\regla{DimensionesOpcionales}{\eps}\nobreak\par
{\leftskip=1.6em\relax\noindent\textup{Primeros} $=$ \{ \eps \}\par}
\par\hangindent=2.2em\hangafter=1
\regla{TipoReferencia}{\tl{idClase} \nt{TipoGenericoOpcional}}\nobreak\par
{\leftskip=1.6em\relax\noindent\textup{Primeros} $=$ \{ \tl{idClase} \}\par}
\par\hangindent=2.2em\hangafter=1
\regla{TipoPrimitivo}{\tl{boolean}}\nobreak\par
{\leftskip=1.6em\relax\noindent\textup{Primeros} $=$ \{ \tl{boolean} \}\par}
\par\hangindent=2.2em\hangafter=1
\regla{TipoPrimitivo}{\tl{char}}\nobreak\par
{\leftskip=1.6em\relax\noindent\textup{Primeros} $=$ \{ \tl{char} \}\par}
\par\hangindent=2.2em\hangafter=1
\regla{TipoPrimitivo}{\tl{int}}\nobreak\par
{\leftskip=1.6em\relax\noindent\textup{Primeros} $=$ \{ \tl{int} \}\par}
\par\hangindent=2.2em\hangafter=1
\regla{TipoPrimitivo}{\tl{float}}\nobreak\par
{\leftskip=1.6em\relax\noindent\textup{Primeros} $=$ \{ \tl{float} \}\par}
\par\hangindent=2.2em\hangafter=1
\regla{TipoGenericoOpcional}{\tl{<} \nt{ArgumentoDeTipo} \tl{>}}\nobreak\par
{\leftskip=1.6em\relax\noindent\textup{Primeros} $=$ \{ \tl{<} \}\par}
\par\hangindent=2.2em\hangafter=1
\regla{TipoGenericoOpcional}{\eps}\nobreak\par
{\leftskip=1.6em\relax\noindent\textup{Primeros} $=$ \{ \eps \}\par}
\par\hangindent=2.2em\hangafter=1
\regla{ArgumentoDeTipo}{\tl{idGen}}\nobreak\par
{\leftskip=1.6em\relax\noindent\textup{Primeros} $=$ \{ \tl{idGen} \}\par}
\par\hangindent=2.2em\hangafter=1
\regla{ArgumentoDeTipo}{\tl{idClase} \nt{TipoGenericoOpcional}}\nobreak\par
{\leftskip=1.6em\relax\noindent\textup{Primeros} $=$ \{ \tl{idClase} \}\par}
\par\hangindent=2.2em\hangafter=1
\regla{ArgumentoDeTipo}{\tl{idClase}}\nobreak\par
{\leftskip=1.6em\relax\noindent\textup{Primeros} $=$ \{ \tl{idClase} \}\par}
\par\hangindent=2.2em\hangafter=1
\regla{ArgsFormales}{\tl{(} \nt{ListaArgsFormalesOpcional} \tl{)}}\nobreak\par
{\leftskip=1.6em\relax\noindent\textup{Primeros} $=$ \{ \tl{(} \}\par}
\par\hangindent=2.2em\hangafter=1
\regla{ListaArgsFormalesOpcional}{\nt{ListaArgsFormales}}\nobreak\par
{\leftskip=1.6em\relax\noindent\textup{Primeros} $=$ \{ \tl{boolean}, \tl{char}, \tl{float}, \tl{idClase}, \tl{idGen}, \tl{int} \}\par}
\par\hangindent=2.2em\hangafter=1
\regla{ListaArgsFormalesOpcional}{\eps}\nobreak\par
{\leftskip=1.6em\relax\noindent\textup{Primeros} $=$ \{ \eps \}\par}
\par\hangindent=2.2em\hangafter=1
\regla{ListaArgsFormales}{\nt{ArgFormal} \nt{RestoListaArgsFormales}}\nobreak\par
{\leftskip=1.6em\relax\noindent\textup{Primeros} $=$ \{ \tl{boolean}, \tl{char}, \tl{float}, \tl{idClase}, \tl{idGen}, \tl{int} \}\par}
\par\hangindent=2.2em\hangafter=1
\regla{RestoListaArgsFormales}{\tl{,} \nt{ArgFormal} \nt{RestoListaArgsFormales}}\nobreak\par
{\leftskip=1.6em\relax\noindent\textup{Primeros} $=$ \{ \tl{,} \}\par}
\par\hangindent=2.2em\hangafter=1
\regla{RestoListaArgsFormales}{\eps}\nobreak\par
{\leftskip=1.6em\relax\noindent\textup{Primeros} $=$ \{ \eps \}\par}
\par\hangindent=2.2em\hangafter=1
\regla{ArgFormal}{\nt{Tipo} \tl{idMetVar}}\nobreak\par
{\leftskip=1.6em\relax\noindent\textup{Primeros} $=$ \{ \tl{boolean}, \tl{char}, \tl{float}, \tl{idClase}, \tl{idGen}, \tl{int} \}\par}
\par\hangindent=2.2em\hangafter=1
\regla{Bloque}{\tl{\{} \nt{ListaSentencias} \tl{\}}}\nobreak\par
{\leftskip=1.6em\relax\noindent\textup{Primeros} $=$ \{ \tl{\{} \}\par}
\par\hangindent=2.2em\hangafter=1
\regla{ListaSentencias}{\nt{Sentencia} \nt{ListaSentencias}}\nobreak\par
{\leftskip=1.6em\relax\noindent\textup{Primeros} $=$ \{ \tl{!}, \tl{(}, \tl{+}, \tl{-}, \tl{;}, \tl{boolean}, \tl{char}, \tl{charLiteral}, \tl{false}, \tl{float}, \tl{floatLiteral}, \tl{for}, \tl{idClase}, \tl{idGen}, \tl{idMetVar}, \tl{if}, \tl{int}, \tl{intLiteral}, \tl{new}, \tl{null}, \tl{return}, \tl{stringLiteral}, \tl{this}, \tl{true}, \tl{var}, \tl{while}, \tl{\{} \}\par}
\par\hangindent=2.2em\hangafter=1
\regla{ListaSentencias}{\eps}\nobreak\par
{\leftskip=1.6em\relax\noindent\textup{Primeros} $=$ \{ \eps \}\par}
\par\hangindent=2.2em\hangafter=1
\regla{Sentencia}{\tl{;}}\nobreak\par
{\leftskip=1.6em\relax\noindent\textup{Primeros} $=$ \{ \tl{;} \}\par}
\par\hangindent=2.2em\hangafter=1
\regla{Sentencia}{\nt{VarLocal} \tl{;}}\nobreak\par
{\leftskip=1.6em\relax\noindent\textup{Primeros} $=$ \{ \tl{var} \}\par}
\par\hangindent=2.2em\hangafter=1
\regla{Sentencia}{\nt{Return} \tl{;}}\nobreak\par
{\leftskip=1.6em\relax\noindent\textup{Primeros} $=$ \{ \tl{return} \}\par}
\par\hangindent=2.2em\hangafter=1
\regla{Sentencia}{\nt{If}}\nobreak\par
{\leftskip=1.6em\relax\noindent\textup{Primeros} $=$ \{ \tl{if} \}\par}
\par\hangindent=2.2em\hangafter=1
\regla{Sentencia}{\nt{While}}\nobreak\par
{\leftskip=1.6em\relax\noindent\textup{Primeros} $=$ \{ \tl{while} \}\par}
\par\hangindent=2.2em\hangafter=1
\regla{Sentencia}{\nt{Bloque}}\nobreak\par
{\leftskip=1.6em\relax\noindent\textup{Primeros} $=$ \{ \tl{\{} \}\par}
\par\hangindent=2.2em\hangafter=1
\regla{Sentencia}{\nt{For}}\nobreak\par
{\leftskip=1.6em\relax\noindent\textup{Primeros} $=$ \{ \tl{for} \}\par}
\par\hangindent=2.2em\hangafter=1
\regla{Sentencia}{\nt{SentenciaExpresion} \tl{;}}\nobreak\par
{\leftskip=1.6em\relax\noindent\textup{Primeros} $=$ \{ \tl{!}, \tl{(}, \tl{+}, \tl{-}, \tl{charLiteral}, \tl{false}, \tl{floatLiteral}, \tl{idClase}, \tl{idMetVar}, \tl{intLiteral}, \tl{new}, \tl{null}, \tl{stringLiteral}, \tl{this}, \tl{true} \}\par}
\par\hangindent=2.2em\hangafter=1
\regla{Sentencia}{\nt{ExpresionSinIdClase} \tl{;}}\nobreak\par
{\leftskip=1.6em\relax\noindent\textup{Primeros} $=$ \{ \tl{!}, \tl{(}, \tl{+}, \tl{-}, \tl{charLiteral}, \tl{false}, \tl{floatLiteral}, \tl{idMetVar}, \tl{intLiteral}, \tl{new}, \tl{null}, \tl{stringLiteral}, \tl{this}, \tl{true} \}\par}
\par\hangindent=2.2em\hangafter=1
\regla{Sentencia}{\nt{TipoPrimitivo} \nt{DimensionesOpcionales} \nt{ListaDeclaradores} \tl{;}}\nobreak\par
{\leftskip=1.6em\relax\noindent\textup{Primeros} $=$ \{ \tl{boolean}, \tl{char}, \tl{float}, \tl{int} \}\par}
\par\hangindent=2.2em\hangafter=1
\regla{Sentencia}{\tl{idGen} \nt{DimensionesOpcionales} \nt{ListaDeclaradores} \tl{;}}\nobreak\par
{\leftskip=1.6em\relax\noindent\textup{Primeros} $=$ \{ \tl{idGen} \}\par}
\par\hangindent=2.2em\hangafter=1
\regla{Sentencia}{\tl{idClase} \nt{SentenciaDesdeIdClase}}\nobreak\par
{\leftskip=1.6em\relax\noindent\textup{Primeros} $=$ \{ \tl{idClase} \}\par}
\par\hangindent=2.2em\hangafter=1
\regla{SentenciaExpresion}{\nt{Expresion}}\nobreak\par
{\leftskip=1.6em\relax\noindent\textup{Primeros} $=$ \{ \tl{!}, \tl{(}, \tl{+}, \tl{-}, \tl{charLiteral}, \tl{false}, \tl{floatLiteral}, \tl{idClase}, \tl{idMetVar}, \tl{intLiteral}, \tl{new}, \tl{null}, \tl{stringLiteral}, \tl{this}, \tl{true} \}\par}
\par\hangindent=2.2em\hangafter=1
\regla{SentenciaDesdeIdClase}{\nt{TipoGenericoOpcional} \nt{DimensionesOpcionales} \nt{ListaDeclaradores} \tl{;}}\nobreak\par
{\leftskip=1.6em\relax\noindent\textup{Primeros} $=$ \{ \tl{<}, \tl{[}, \tl{idMetVar} \}\par}
\par\hangindent=2.2em\hangafter=1
\regla{SentenciaDesdeIdClase}{\tl{.} \tl{idMetVar} \nt{ArgsActuales} \nt{ColaDeExpresionDesdePrimario} \tl{;}}\nobreak\par
{\leftskip=1.6em\relax\noindent\textup{Primeros} $=$ \{ \tl{.} \}\par}
\par\hangindent=2.2em\hangafter=1
\regla{VarLocal}{\tl{var} \tl{idMetVar} \tl{=} \nt{ExpresionCondicional}}\nobreak\par
{\leftskip=1.6em\relax\noindent\textup{Primeros} $=$ \{ \tl{var} \}\par}
\par\hangindent=2.2em\hangafter=1
\regla{ListaDeclaradores}{\tl{idMetVar} \nt{InicializacionDeDeclarador} \nt{RestoListaDeclaradores}}\nobreak\par
{\leftskip=1.6em\relax\noindent\textup{Primeros} $=$ \{ \tl{idMetVar} \}\par}
\par\hangindent=2.2em\hangafter=1
\regla{RestoListaDeclaradores}{\tl{,} \tl{idMetVar} \nt{InicializacionDeDeclarador} \nt{RestoListaDeclaradores}}\nobreak\par
{\leftskip=1.6em\relax\noindent\textup{Primeros} $=$ \{ \tl{,} \}\par}
\par\hangindent=2.2em\hangafter=1
\regla{RestoListaDeclaradores}{\eps}\nobreak\par
{\leftskip=1.6em\relax\noindent\textup{Primeros} $=$ \{ \eps \}\par}
\par\hangindent=2.2em\hangafter=1
\regla{InicializacionDeDeclarador}{\tl{=} \nt{ExpresionCondicional}}\nobreak\par
{\leftskip=1.6em\relax\noindent\textup{Primeros} $=$ \{ \tl{=} \}\par}
\par\hangindent=2.2em\hangafter=1
\regla{InicializacionDeDeclarador}{\eps}\nobreak\par
{\leftskip=1.6em\relax\noindent\textup{Primeros} $=$ \{ \eps \}\par}
\par\hangindent=2.2em\hangafter=1
\regla{Return}{\tl{return} \nt{ExpresionOpcional}}\nobreak\par
{\leftskip=1.6em\relax\noindent\textup{Primeros} $=$ \{ \tl{return} \}\par}
\par\hangindent=2.2em\hangafter=1
\regla{ExpresionOpcional}{\nt{Expresion}}\nobreak\par
{\leftskip=1.6em\relax\noindent\textup{Primeros} $=$ \{ \tl{!}, \tl{(}, \tl{+}, \tl{-}, \tl{charLiteral}, \tl{false}, \tl{floatLiteral}, \tl{idClase}, \tl{idMetVar}, \tl{intLiteral}, \tl{new}, \tl{null}, \tl{stringLiteral}, \tl{this}, \tl{true} \}\par}
\par\hangindent=2.2em\hangafter=1
\regla{ExpresionOpcional}{\eps}\nobreak\par
{\leftskip=1.6em\relax\noindent\textup{Primeros} $=$ \{ \eps \}\par}
\par\hangindent=2.2em\hangafter=1
\regla{If}{\tl{if} \tl{(} \nt{Expresion} \tl{)} \nt{Sentencia} \nt{ElseOpcional}}\nobreak\par
{\leftskip=1.6em\relax\noindent\textup{Primeros} $=$ \{ \tl{if} \}\par}
\par\hangindent=2.2em\hangafter=1
\regla{ElseOpcional}{\tl{else} \nt{Sentencia}}\nobreak\par
{\leftskip=1.6em\relax\noindent\textup{Primeros} $=$ \{ \tl{else} \}\par}
\par\hangindent=2.2em\hangafter=1
\regla{ElseOpcional}{\eps}\nobreak\par
{\leftskip=1.6em\relax\noindent\textup{Primeros} $=$ \{ \eps \}\par}
\par\hangindent=2.2em\hangafter=1
\regla{While}{\tl{while} \tl{(} \nt{Expresion} \tl{)} \nt{Sentencia}}\nobreak\par
{\leftskip=1.6em\relax\noindent\textup{Primeros} $=$ \{ \tl{while} \}\par}
\par\hangindent=2.2em\hangafter=1
\regla{For}{\tl{for} \tl{(} \nt{EncabezadoFor} \tl{)} \nt{Sentencia}}\nobreak\par
{\leftskip=1.6em\relax\noindent\textup{Primeros} $=$ \{ \tl{for} \}\par}
\par\hangindent=2.2em\hangafter=1
\regla{EncabezadoFor}{\tl{;} \nt{ExpresionOpcional} \tl{;} \nt{ExpresionOpcional}}\nobreak\par
{\leftskip=1.6em\relax\noindent\textup{Primeros} $=$ \{ \tl{;} \}\par}
\par\hangindent=2.2em\hangafter=1
\regla{EncabezadoFor}{\tl{var} \tl{idMetVar} \nt{EncabezadoForTrasVar}}\nobreak\par
{\leftskip=1.6em\relax\noindent\textup{Primeros} $=$ \{ \tl{var} \}\par}
\par\hangindent=2.2em\hangafter=1
\regla{EncabezadoFor}{\nt{TipoPrimitivo} \nt{DimensionesOpcionales} \tl{idMetVar} \nt{EncabezadoForTrasDeclarador}}\nobreak\par
{\leftskip=1.6em\relax\noindent\textup{Primeros} $=$ \{ \tl{boolean}, \tl{char}, \tl{float}, \tl{int} \}\par}
\par\hangindent=2.2em\hangafter=1
\regla{EncabezadoFor}{\tl{idGen} \nt{DimensionesOpcionales} \tl{idMetVar} \nt{EncabezadoForTrasDeclarador}}\nobreak\par
{\leftskip=1.6em\relax\noindent\textup{Primeros} $=$ \{ \tl{idGen} \}\par}
\par\hangindent=2.2em\hangafter=1
\regla{EncabezadoFor}{\tl{idClase} \nt{EncabezadoForDesdeIdClase}}\nobreak\par
{\leftskip=1.6em\relax\noindent\textup{Primeros} $=$ \{ \tl{idClase} \}\par}
\par\hangindent=2.2em\hangafter=1
\regla{EncabezadoFor}{\nt{ExpresionSinIdClase} \nt{ClausulasFinalesFor}}\nobreak\par
{\leftskip=1.6em\relax\noindent\textup{Primeros} $=$ \{ \tl{!}, \tl{(}, \tl{+}, \tl{-}, \tl{charLiteral}, \tl{false}, \tl{floatLiteral}, \tl{idMetVar}, \tl{intLiteral}, \tl{new}, \tl{null}, \tl{stringLiteral}, \tl{this}, \tl{true} \}\par}
\par\hangindent=2.2em\hangafter=1
\regla{EncabezadoForDesdeIdClase}{\nt{TipoGenericoOpcional} \nt{DimensionesOpcionales} \tl{idMetVar} \nt{EncabezadoForTrasDeclarador}}\nobreak\par
{\leftskip=1.6em\relax\noindent\textup{Primeros} $=$ \{ \tl{<}, \tl{[}, \tl{idMetVar} \}\par}
\par\hangindent=2.2em\hangafter=1
\regla{EncabezadoForDesdeIdClase}{\tl{.} \tl{idMetVar} \nt{ArgsActuales} \nt{ColaDeExpresionDesdePrimario} \nt{ClausulasFinalesFor}}\nobreak\par
{\leftskip=1.6em\relax\noindent\textup{Primeros} $=$ \{ \tl{.} \}\par}
\par\hangindent=2.2em\hangafter=1
\regla{EncabezadoForTrasVar}{\tl{:} \nt{Expresion}}\nobreak\par
{\leftskip=1.6em\relax\noindent\textup{Primeros} $=$ \{ \tl{:} \}\par}
\par\hangindent=2.2em\hangafter=1
\regla{EncabezadoForTrasVar}{\tl{=} \nt{ExpresionCondicional} \nt{ClausulasFinalesFor}}\nobreak\par
{\leftskip=1.6em\relax\noindent\textup{Primeros} $=$ \{ \tl{=} \}\par}
\par\hangindent=2.2em\hangafter=1
\regla{EncabezadoForTrasDeclarador}{\tl{:} \nt{Expresion}}\nobreak\par
{\leftskip=1.6em\relax\noindent\textup{Primeros} $=$ \{ \tl{:} \}\par}
\par\hangindent=2.2em\hangafter=1
\regla{EncabezadoForTrasDeclarador}{\nt{InicializacionDeDeclaradorFor} \nt{ClausulasFinalesFor}}\nobreak\par
{\leftskip=1.6em\relax\noindent\textup{Primeros} $=$ \{ \tl{;}, \tl{=} \}\par}
\par\hangindent=2.2em\hangafter=1
\regla{InicializacionDeDeclaradorFor}{\tl{=} \nt{ExpresionCondicional}}\nobreak\par
{\leftskip=1.6em\relax\noindent\textup{Primeros} $=$ \{ \tl{=} \}\par}
\par\hangindent=2.2em\hangafter=1
\regla{InicializacionDeDeclaradorFor}{\eps}\nobreak\par
{\leftskip=1.6em\relax\noindent\textup{Primeros} $=$ \{ \eps \}\par}
\par\hangindent=2.2em\hangafter=1
\regla{ClausulasFinalesFor}{\tl{;} \nt{ExpresionOpcional} \tl{;} \nt{ExpresionOpcional}}\nobreak\par
{\leftskip=1.6em\relax\noindent\textup{Primeros} $=$ \{ \tl{;} \}\par}
\par\hangindent=2.2em\hangafter=1
\regla{Expresion}{\nt{ExpresionCondicional} \nt{RestoExpresion}}\nobreak\par
{\leftskip=1.6em\relax\noindent\textup{Primeros} $=$ \{ \tl{!}, \tl{(}, \tl{+}, \tl{-}, \tl{charLiteral}, \tl{false}, \tl{floatLiteral}, \tl{idClase}, \tl{idMetVar}, \tl{intLiteral}, \tl{new}, \tl{null}, \tl{stringLiteral}, \tl{this}, \tl{true} \}\par}
\par\hangindent=2.2em\hangafter=1
\regla{RestoExpresion}{\nt{OperadorAsignacion} \nt{ExpresionCondicional}}\nobreak\par
{\leftskip=1.6em\relax\noindent\textup{Primeros} $=$ \{ \tl{=} \}\par}
\par\hangindent=2.2em\hangafter=1
\regla{RestoExpresion}{\eps}\nobreak\par
{\leftskip=1.6em\relax\noindent\textup{Primeros} $=$ \{ \eps \}\par}
\par\hangindent=2.2em\hangafter=1
\regla{OperadorAsignacion}{\tl{=}}\nobreak\par
{\leftskip=1.6em\relax\noindent\textup{Primeros} $=$ \{ \tl{=} \}\par}
\par\hangindent=2.2em\hangafter=1
\regla{ExpresionCondicional}{\nt{ExpresionCompuesta} \nt{RestoCondicional}}\nobreak\par
{\leftskip=1.6em\relax\noindent\textup{Primeros} $=$ \{ \tl{!}, \tl{(}, \tl{+}, \tl{-}, \tl{charLiteral}, \tl{false}, \tl{floatLiteral}, \tl{idClase}, \tl{idMetVar}, \tl{intLiteral}, \tl{new}, \tl{null}, \tl{stringLiteral}, \tl{this}, \tl{true} \}\par}
\par\hangindent=2.2em\hangafter=1
\regla{RestoCondicional}{\tl{?} \nt{Expresion} \tl{:} \nt{ExpresionCondicional}}\nobreak\par
{\leftskip=1.6em\relax\noindent\textup{Primeros} $=$ \{ \tl{?} \}\par}
\par\hangindent=2.2em\hangafter=1
\regla{RestoCondicional}{\eps}\nobreak\par
{\leftskip=1.6em\relax\noindent\textup{Primeros} $=$ \{ \eps \}\par}
\par\hangindent=2.2em\hangafter=1
\regla{ExpresionCompuesta}{\nt{ExpresionBasica} \nt{RestoExpresionCompuesta}}\nobreak\par
{\leftskip=1.6em\relax\noindent\textup{Primeros} $=$ \{ \tl{!}, \tl{(}, \tl{+}, \tl{-}, \tl{charLiteral}, \tl{false}, \tl{floatLiteral}, \tl{idClase}, \tl{idMetVar}, \tl{intLiteral}, \tl{new}, \tl{null}, \tl{stringLiteral}, \tl{this}, \tl{true} \}\par}
\par\hangindent=2.2em\hangafter=1
\regla{RestoExpresionCompuesta}{\nt{OperadorBinario} \nt{ExpresionBasica} \nt{RestoExpresionCompuesta}}\nobreak\par
{\leftskip=1.6em\relax\noindent\textup{Primeros} $=$ \{ \tl{!=}, \tl{\%}, \tl{\&\&}, \tl{*}, \tl{+}, \tl{-}, \tl{/}, \tl{<}, \tl{<=}, \tl{==}, \tl{>}, \tl{>=}, \tl{||} \}\par}
\par\hangindent=2.2em\hangafter=1
\regla{RestoExpresionCompuesta}{\eps}\nobreak\par
{\leftskip=1.6em\relax\noindent\textup{Primeros} $=$ \{ \eps \}\par}
\par\hangindent=2.2em\hangafter=1
\regla{OperadorBinario}{\tl{||}}\nobreak\par
{\leftskip=1.6em\relax\noindent\textup{Primeros} $=$ \{ \tl{||} \}\par}
\par\hangindent=2.2em\hangafter=1
\regla{OperadorBinario}{\tl{\&\&}}\nobreak\par
{\leftskip=1.6em\relax\noindent\textup{Primeros} $=$ \{ \tl{\&\&} \}\par}
\par\hangindent=2.2em\hangafter=1
\regla{OperadorBinario}{\tl{==}}\nobreak\par
{\leftskip=1.6em\relax\noindent\textup{Primeros} $=$ \{ \tl{==} \}\par}
\par\hangindent=2.2em\hangafter=1
\regla{OperadorBinario}{\tl{!=}}\nobreak\par
{\leftskip=1.6em\relax\noindent\textup{Primeros} $=$ \{ \tl{!=} \}\par}
\par\hangindent=2.2em\hangafter=1
\regla{OperadorBinario}{\tl{<}}\nobreak\par
{\leftskip=1.6em\relax\noindent\textup{Primeros} $=$ \{ \tl{<} \}\par}
\par\hangindent=2.2em\hangafter=1
\regla{OperadorBinario}{\tl{>}}\nobreak\par
{\leftskip=1.6em\relax\noindent\textup{Primeros} $=$ \{ \tl{>} \}\par}
\par\hangindent=2.2em\hangafter=1
\regla{OperadorBinario}{\tl{<=}}\nobreak\par
{\leftskip=1.6em\relax\noindent\textup{Primeros} $=$ \{ \tl{<=} \}\par}
\par\hangindent=2.2em\hangafter=1
\regla{OperadorBinario}{\tl{>=}}\nobreak\par
{\leftskip=1.6em\relax\noindent\textup{Primeros} $=$ \{ \tl{>=} \}\par}
\par\hangindent=2.2em\hangafter=1
\regla{OperadorBinario}{\tl{+}}\nobreak\par
{\leftskip=1.6em\relax\noindent\textup{Primeros} $=$ \{ \tl{+} \}\par}
\par\hangindent=2.2em\hangafter=1
\regla{OperadorBinario}{\tl{-}}\nobreak\par
{\leftskip=1.6em\relax\noindent\textup{Primeros} $=$ \{ \tl{-} \}\par}
\par\hangindent=2.2em\hangafter=1
\regla{OperadorBinario}{\tl{*}}\nobreak\par
{\leftskip=1.6em\relax\noindent\textup{Primeros} $=$ \{ \tl{*} \}\par}
\par\hangindent=2.2em\hangafter=1
\regla{OperadorBinario}{\tl{/}}\nobreak\par
{\leftskip=1.6em\relax\noindent\textup{Primeros} $=$ \{ \tl{/} \}\par}
\par\hangindent=2.2em\hangafter=1
\regla{OperadorBinario}{\tl{\%}}\nobreak\par
{\leftskip=1.6em\relax\noindent\textup{Primeros} $=$ \{ \tl{\%} \}\par}
\par\hangindent=2.2em\hangafter=1
\regla{ExpresionBasica}{\nt{OperadorUnario} \nt{Operando}}\nobreak\par
{\leftskip=1.6em\relax\noindent\textup{Primeros} $=$ \{ \tl{!}, \tl{+}, \tl{-} \}\par}
\par\hangindent=2.2em\hangafter=1
\regla{ExpresionBasica}{\nt{Operando}}\nobreak\par
{\leftskip=1.6em\relax\noindent\textup{Primeros} $=$ \{ \tl{(}, \tl{charLiteral}, \tl{false}, \tl{floatLiteral}, \tl{idClase}, \tl{idMetVar}, \tl{intLiteral}, \tl{new}, \tl{null}, \tl{stringLiteral}, \tl{this}, \tl{true} \}\par}
\par\hangindent=2.2em\hangafter=1
\regla{OperadorUnario}{\tl{+}}\nobreak\par
{\leftskip=1.6em\relax\noindent\textup{Primeros} $=$ \{ \tl{+} \}\par}
\par\hangindent=2.2em\hangafter=1
\regla{OperadorUnario}{\tl{-}}\nobreak\par
{\leftskip=1.6em\relax\noindent\textup{Primeros} $=$ \{ \tl{-} \}\par}
\par\hangindent=2.2em\hangafter=1
\regla{OperadorUnario}{\tl{!}}\nobreak\par
{\leftskip=1.6em\relax\noindent\textup{Primeros} $=$ \{ \tl{!} \}\par}
\par\hangindent=2.2em\hangafter=1
\regla{Operando}{\nt{Primitivo}}\nobreak\par
{\leftskip=1.6em\relax\noindent\textup{Primeros} $=$ \{ \tl{charLiteral}, \tl{false}, \tl{floatLiteral}, \tl{intLiteral}, \tl{null}, \tl{true} \}\par}
\par\hangindent=2.2em\hangafter=1
\regla{Operando}{\nt{Referencia} \nt{PosfijoOpcional}}\nobreak\par
{\leftskip=1.6em\relax\noindent\textup{Primeros} $=$ \{ \tl{(}, \tl{idClase}, \tl{idMetVar}, \tl{new}, \tl{stringLiteral}, \tl{this} \}\par}
\par\hangindent=2.2em\hangafter=1
\regla{PosfijoOpcional}{\tl{++}}\nobreak\par
{\leftskip=1.6em\relax\noindent\textup{Primeros} $=$ \{ \tl{++} \}\par}
\par\hangindent=2.2em\hangafter=1
\regla{PosfijoOpcional}{\tl{-{}-}}\nobreak\par
{\leftskip=1.6em\relax\noindent\textup{Primeros} $=$ \{ \tl{-{}-} \}\par}
\par\hangindent=2.2em\hangafter=1
\regla{PosfijoOpcional}{\eps}\nobreak\par
{\leftskip=1.6em\relax\noindent\textup{Primeros} $=$ \{ \eps \}\par}
\par\hangindent=2.2em\hangafter=1
\regla{Primitivo}{\tl{true}}\nobreak\par
{\leftskip=1.6em\relax\noindent\textup{Primeros} $=$ \{ \tl{true} \}\par}
\par\hangindent=2.2em\hangafter=1
\regla{Primitivo}{\tl{false}}\nobreak\par
{\leftskip=1.6em\relax\noindent\textup{Primeros} $=$ \{ \tl{false} \}\par}
\par\hangindent=2.2em\hangafter=1
\regla{Primitivo}{\tl{intLiteral}}\nobreak\par
{\leftskip=1.6em\relax\noindent\textup{Primeros} $=$ \{ \tl{intLiteral} \}\par}
\par\hangindent=2.2em\hangafter=1
\regla{Primitivo}{\tl{charLiteral}}\nobreak\par
{\leftskip=1.6em\relax\noindent\textup{Primeros} $=$ \{ \tl{charLiteral} \}\par}
\par\hangindent=2.2em\hangafter=1
\regla{Primitivo}{\tl{null}}\nobreak\par
{\leftskip=1.6em\relax\noindent\textup{Primeros} $=$ \{ \tl{null} \}\par}
\par\hangindent=2.2em\hangafter=1
\regla{Primitivo}{\tl{floatLiteral}}\nobreak\par
{\leftskip=1.6em\relax\noindent\textup{Primeros} $=$ \{ \tl{floatLiteral} \}\par}
\par\hangindent=2.2em\hangafter=1
\regla{ExpresionSinIdClase}{\nt{OperadorUnario} \nt{Operando} \nt{ColaTrasOperando}}\nobreak\par
{\leftskip=1.6em\relax\noindent\textup{Primeros} $=$ \{ \tl{!}, \tl{+}, \tl{-} \}\par}
\par\hangindent=2.2em\hangafter=1
\regla{ExpresionSinIdClase}{\nt{Primitivo} \nt{ColaTrasOperando}}\nobreak\par
{\leftskip=1.6em\relax\noindent\textup{Primeros} $=$ \{ \tl{charLiteral}, \tl{false}, \tl{floatLiteral}, \tl{intLiteral}, \tl{null}, \tl{true} \}\par}
\par\hangindent=2.2em\hangafter=1
\regla{ExpresionSinIdClase}{\nt{PrimarioSinIdClase} \nt{ColaDeExpresionDesdePrimario}}\nobreak\par
{\leftskip=1.6em\relax\noindent\textup{Primeros} $=$ \{ \tl{(}, \tl{idMetVar}, \tl{new}, \tl{stringLiteral}, \tl{this} \}\par}
\par\hangindent=2.2em\hangafter=1
\regla{ColaTrasOperando}{\nt{RestoExpresionCompuesta} \nt{RestoCondicional} \nt{RestoExpresion}}\nobreak\par
{\leftskip=1.6em\relax\noindent\textup{Primeros} $=$ \{ \tl{!=}, \tl{\%}, \tl{\&\&}, \tl{*}, \tl{+}, \tl{-}, \tl{/}, \tl{<}, \tl{<=}, \tl{=}, \tl{==}, \tl{>}, \tl{>=}, \tl{?}, \tl{||}, \eps \}\par}
\par\hangindent=2.2em\hangafter=1
\regla{ColaDeExpresionDesdePrimario}{\nt{RestoReferencia} \nt{PosfijoOpcional} \nt{ColaTrasOperando}}\nobreak\par
{\leftskip=1.6em\relax\noindent\textup{Primeros} $=$ \{ \tl{!=}, \tl{\%}, \tl{\&\&}, \tl{*}, \tl{+}, \tl{++}, \tl{-}, \tl{-{}-}, \tl{.}, \tl{/}, \tl{<}, \tl{<=}, \tl{=}, \tl{==}, \tl{>}, \tl{>=}, \tl{?}, \tl{[}, \tl{||}, \eps \}\par}
\par\hangindent=2.2em\hangafter=1
\regla{Referencia}{\nt{Primario} \nt{RestoReferencia}}\nobreak\par
{\leftskip=1.6em\relax\noindent\textup{Primeros} $=$ \{ \tl{(}, \tl{idClase}, \tl{idMetVar}, \tl{new}, \tl{stringLiteral}, \tl{this} \}\par}
\par\hangindent=2.2em\hangafter=1
\regla{RestoReferencia}{\nt{Encadenado} \nt{RestoReferencia}}\nobreak\par
{\leftskip=1.6em\relax\noindent\textup{Primeros} $=$ \{ \tl{.} \}\par}
\par\hangindent=2.2em\hangafter=1
\regla{RestoReferencia}{\nt{AccesoArreglo} \nt{RestoReferencia}}\nobreak\par
{\leftskip=1.6em\relax\noindent\textup{Primeros} $=$ \{ \tl{[} \}\par}
\par\hangindent=2.2em\hangafter=1
\regla{RestoReferencia}{\eps}\nobreak\par
{\leftskip=1.6em\relax\noindent\textup{Primeros} $=$ \{ \eps \}\par}
\par\hangindent=2.2em\hangafter=1
\regla{Encadenado}{\tl{.} \tl{idMetVar} \nt{EncadenadoAtributoOMetodo}}\nobreak\par
{\leftskip=1.6em\relax\noindent\textup{Primeros} $=$ \{ \tl{.} \}\par}
\par\hangindent=2.2em\hangafter=1
\regla{EncadenadoAtributoOMetodo}{\nt{ArgsActuales}}\nobreak\par
{\leftskip=1.6em\relax\noindent\textup{Primeros} $=$ \{ \tl{(} \}\par}
\par\hangindent=2.2em\hangafter=1
\regla{EncadenadoAtributoOMetodo}{\eps}\nobreak\par
{\leftskip=1.6em\relax\noindent\textup{Primeros} $=$ \{ \eps \}\par}
\par\hangindent=2.2em\hangafter=1
\regla{AccesoArreglo}{\tl{[} \nt{Expresion} \tl{]}}\nobreak\par
{\leftskip=1.6em\relax\noindent\textup{Primeros} $=$ \{ \tl{[} \}\par}
\par\hangindent=2.2em\hangafter=1
\regla{Primario}{\nt{LlamadaMetodoEstatico}}\nobreak\par
{\leftskip=1.6em\relax\noindent\textup{Primeros} $=$ \{ \tl{idClase} \}\par}
\par\hangindent=2.2em\hangafter=1
\regla{Primario}{\nt{PrimarioSinIdClase}}\nobreak\par
{\leftskip=1.6em\relax\noindent\textup{Primeros} $=$ \{ \tl{(}, \tl{idMetVar}, \tl{new}, \tl{stringLiteral}, \tl{this} \}\par}
\par\hangindent=2.2em\hangafter=1
\regla{LlamadaMetodoEstatico}{\tl{idClase} \tl{.} \tl{idMetVar} \nt{ArgsActuales}}\nobreak\par
{\leftskip=1.6em\relax\noindent\textup{Primeros} $=$ \{ \tl{idClase} \}\par}
\par\hangindent=2.2em\hangafter=1
\regla{PrimarioSinIdClase}{\tl{this}}\nobreak\par
{\leftskip=1.6em\relax\noindent\textup{Primeros} $=$ \{ \tl{this} \}\par}
\par\hangindent=2.2em\hangafter=1
\regla{PrimarioSinIdClase}{\tl{stringLiteral}}\nobreak\par
{\leftskip=1.6em\relax\noindent\textup{Primeros} $=$ \{ \tl{stringLiteral} \}\par}
\par\hangindent=2.2em\hangafter=1
\regla{PrimarioSinIdClase}{\nt{ParentizadaOLambda}}\nobreak\par
{\leftskip=1.6em\relax\noindent\textup{Primeros} $=$ \{ \tl{(} \}\par}
\par\hangindent=2.2em\hangafter=1
\regla{PrimarioSinIdClase}{\nt{AccesoVarLlamadaOLambda}}\nobreak\par
{\leftskip=1.6em\relax\noindent\textup{Primeros} $=$ \{ \tl{idMetVar} \}\par}
\par\hangindent=2.2em\hangafter=1
\regla{PrimarioSinIdClase}{\nt{Creacion}}\nobreak\par
{\leftskip=1.6em\relax\noindent\textup{Primeros} $=$ \{ \tl{new} \}\par}
\par\hangindent=2.2em\hangafter=1
\regla{AccesoVarLlamadaOLambda}{\tl{idMetVar} \nt{AccesoLlamadaOLambda}}\nobreak\par
{\leftskip=1.6em\relax\noindent\textup{Primeros} $=$ \{ \tl{idMetVar} \}\par}
\par\hangindent=2.2em\hangafter=1
\regla{AccesoLlamadaOLambda}{\nt{ArgsActuales}}\nobreak\par
{\leftskip=1.6em\relax\noindent\textup{Primeros} $=$ \{ \tl{(} \}\par}
\par\hangindent=2.2em\hangafter=1
\regla{AccesoLlamadaOLambda}{\tl{->} \nt{CuerpoLambda}}\nobreak\par
{\leftskip=1.6em\relax\noindent\textup{Primeros} $=$ \{ \tl{->} \}\par}
\par\hangindent=2.2em\hangafter=1
\regla{AccesoLlamadaOLambda}{\eps}\nobreak\par
{\leftskip=1.6em\relax\noindent\textup{Primeros} $=$ \{ \eps \}\par}
\par\hangindent=2.2em\hangafter=1
\regla{ParentizadaOLambda}{\tl{(} \nt{ContenidoDeParentesis}}\nobreak\par
{\leftskip=1.6em\relax\noindent\textup{Primeros} $=$ \{ \tl{(} \}\par}
\par\hangindent=2.2em\hangafter=1
\regla{ContenidoDeParentesis}{\tl{)} \tl{->} \nt{CuerpoLambda}}\nobreak\par
{\leftskip=1.6em\relax\noindent\textup{Primeros} $=$ \{ \tl{)} \}\par}
\par\hangindent=2.2em\hangafter=1
\regla{ContenidoDeParentesis}{\nt{Expresion} \nt{RestoDeParentesis}}\nobreak\par
{\leftskip=1.6em\relax\noindent\textup{Primeros} $=$ \{ \tl{!}, \tl{(}, \tl{+}, \tl{-}, \tl{charLiteral}, \tl{false}, \tl{floatLiteral}, \tl{idClase}, \tl{idMetVar}, \tl{intLiteral}, \tl{new}, \tl{null}, \tl{stringLiteral}, \tl{this}, \tl{true} \}\par}
\par\hangindent=2.2em\hangafter=1
\regla{RestoDeParentesis}{\tl{)} \nt{FinDeParentizadaOLambda}}\nobreak\par
{\leftskip=1.6em\relax\noindent\textup{Primeros} $=$ \{ \tl{)} \}\par}
\par\hangindent=2.2em\hangafter=1
\regla{RestoDeParentesis}{\tl{,} \nt{ListaParametrosLambda} \tl{)} \tl{->} \nt{CuerpoLambda}}\nobreak\par
{\leftskip=1.6em\relax\noindent\textup{Primeros} $=$ \{ \tl{,} \}\par}
\par\hangindent=2.2em\hangafter=1
\regla{FinDeParentizadaOLambda}{\tl{->} \nt{CuerpoLambda}}\nobreak\par
{\leftskip=1.6em\relax\noindent\textup{Primeros} $=$ \{ \tl{->} \}\par}
\par\hangindent=2.2em\hangafter=1
\regla{FinDeParentizadaOLambda}{\eps}\nobreak\par
{\leftskip=1.6em\relax\noindent\textup{Primeros} $=$ \{ \eps \}\par}
\par\hangindent=2.2em\hangafter=1
\regla{ListaParametrosLambda}{\tl{idMetVar} \nt{RestoListaParametrosLambda}}\nobreak\par
{\leftskip=1.6em\relax\noindent\textup{Primeros} $=$ \{ \tl{idMetVar} \}\par}
\par\hangindent=2.2em\hangafter=1
\regla{RestoListaParametrosLambda}{\tl{,} \tl{idMetVar} \nt{RestoListaParametrosLambda}}\nobreak\par
{\leftskip=1.6em\relax\noindent\textup{Primeros} $=$ \{ \tl{,} \}\par}
\par\hangindent=2.2em\hangafter=1
\regla{RestoListaParametrosLambda}{\eps}\nobreak\par
{\leftskip=1.6em\relax\noindent\textup{Primeros} $=$ \{ \eps \}\par}
\par\hangindent=2.2em\hangafter=1
\regla{CuerpoLambda}{\nt{Expresion}}\nobreak\par
{\leftskip=1.6em\relax\noindent\textup{Primeros} $=$ \{ \tl{!}, \tl{(}, \tl{+}, \tl{-}, \tl{charLiteral}, \tl{false}, \tl{floatLiteral}, \tl{idClase}, \tl{idMetVar}, \tl{intLiteral}, \tl{new}, \tl{null}, \tl{stringLiteral}, \tl{this}, \tl{true} \}\par}
\par\hangindent=2.2em\hangafter=1
\regla{Creacion}{\tl{new} \nt{ArregloOInstancia}}\nobreak\par
{\leftskip=1.6em\relax\noindent\textup{Primeros} $=$ \{ \tl{new} \}\par}
\par\hangindent=2.2em\hangafter=1
\regla{ArregloOInstancia}{\nt{TipoPrimitivo} \nt{CreacionDeArreglo}}\nobreak\par
{\leftskip=1.6em\relax\noindent\textup{Primeros} $=$ \{ \tl{boolean}, \tl{char}, \tl{float}, \tl{int} \}\par}
\par\hangindent=2.2em\hangafter=1
\regla{ArregloOInstancia}{\tl{idGen} \nt{CreacionDeArreglo}}\nobreak\par
{\leftskip=1.6em\relax\noindent\textup{Primeros} $=$ \{ \tl{idGen} \}\par}
\par\hangindent=2.2em\hangafter=1
\regla{ArregloOInstancia}{\tl{idClase} \nt{CreacionDesdeIdClase}}\nobreak\par
{\leftskip=1.6em\relax\noindent\textup{Primeros} $=$ \{ \tl{idClase} \}\par}
\par\hangindent=2.2em\hangafter=1
\regla{CreacionDesdeIdClase}{\tl{<} \nt{GenericoDeCreacion}}\nobreak\par
{\leftskip=1.6em\relax\noindent\textup{Primeros} $=$ \{ \tl{<} \}\par}
\par\hangindent=2.2em\hangafter=1
\regla{CreacionDesdeIdClase}{\nt{CreacionDeArreglo}}\nobreak\par
{\leftskip=1.6em\relax\noindent\textup{Primeros} $=$ \{ \tl{[} \}\par}
\par\hangindent=2.2em\hangafter=1
\regla{CreacionDesdeIdClase}{\nt{ArgsActuales}}\nobreak\par
{\leftskip=1.6em\relax\noindent\textup{Primeros} $=$ \{ \tl{(} \}\par}
\par\hangindent=2.2em\hangafter=1
\regla{GenericoDeCreacion}{\tl{>} \nt{ArgsActuales}}\nobreak\par
{\leftskip=1.6em\relax\noindent\textup{Primeros} $=$ \{ \tl{>} \}\par}
\par\hangindent=2.2em\hangafter=1
\regla{GenericoDeCreacion}{\nt{ArgumentoDeTipo} \tl{>} \nt{ArregloOInstanciaDeClase}}\nobreak\par
{\leftskip=1.6em\relax\noindent\textup{Primeros} $=$ \{ \tl{idClase}, \tl{idGen} \}\par}
\par\hangindent=2.2em\hangafter=1
\regla{ArregloOInstanciaDeClase}{\nt{CreacionDeArreglo}}\nobreak\par
{\leftskip=1.6em\relax\noindent\textup{Primeros} $=$ \{ \tl{[} \}\par}
\par\hangindent=2.2em\hangafter=1
\regla{ArregloOInstanciaDeClase}{\nt{ArgsActuales}}\nobreak\par
{\leftskip=1.6em\relax\noindent\textup{Primeros} $=$ \{ \tl{(} \}\par}
\par\hangindent=2.2em\hangafter=1
\regla{CreacionDeArreglo}{\tl{[} \nt{RestoCreacionDeArreglo}}\nobreak\par
{\leftskip=1.6em\relax\noindent\textup{Primeros} $=$ \{ \tl{[} \}\par}
\par\hangindent=2.2em\hangafter=1
\regla{RestoCreacionDeArreglo}{\nt{Expresion} \tl{]} \nt{DimensionesConTamanioOpcional}}\nobreak\par
{\leftskip=1.6em\relax\noindent\textup{Primeros} $=$ \{ \tl{!}, \tl{(}, \tl{+}, \tl{-}, \tl{charLiteral}, \tl{false}, \tl{floatLiteral}, \tl{idClase}, \tl{idMetVar}, \tl{intLiteral}, \tl{new}, \tl{null}, \tl{stringLiteral}, \tl{this}, \tl{true} \}\par}
\par\hangindent=2.2em\hangafter=1
\regla{RestoCreacionDeArreglo}{\tl{]} \nt{DimensionesVaciasOpcional} \nt{InicializadorDeArreglo}}\nobreak\par
{\leftskip=1.6em\relax\noindent\textup{Primeros} $=$ \{ \tl{]} \}\par}
\par\hangindent=2.2em\hangafter=1
\regla{DimensionesVaciasOpcional}{\tl{[} \tl{]} \nt{DimensionesVaciasOpcional}}\nobreak\par
{\leftskip=1.6em\relax\noindent\textup{Primeros} $=$ \{ \tl{[} \}\par}
\par\hangindent=2.2em\hangafter=1
\regla{DimensionesVaciasOpcional}{\eps}\nobreak\par
{\leftskip=1.6em\relax\noindent\textup{Primeros} $=$ \{ \eps \}\par}
\par\hangindent=2.2em\hangafter=1
\regla{InicializadorDeArreglo}{\tl{\{} \nt{ListaElementosOpcional} \tl{\}}}\nobreak\par
{\leftskip=1.6em\relax\noindent\textup{Primeros} $=$ \{ \tl{\{} \}\par}
\par\hangindent=2.2em\hangafter=1
\regla{ListaElementosOpcional}{\nt{ListaElementos}}\nobreak\par
{\leftskip=1.6em\relax\noindent\textup{Primeros} $=$ \{ \tl{!}, \tl{(}, \tl{+}, \tl{-}, \tl{charLiteral}, \tl{false}, \tl{floatLiteral}, \tl{idClase}, \tl{idMetVar}, \tl{intLiteral}, \tl{new}, \tl{null}, \tl{stringLiteral}, \tl{this}, \tl{true}, \tl{\{} \}\par}
\par\hangindent=2.2em\hangafter=1
\regla{ListaElementosOpcional}{\eps}\nobreak\par
{\leftskip=1.6em\relax\noindent\textup{Primeros} $=$ \{ \eps \}\par}
\par\hangindent=2.2em\hangafter=1
\regla{ListaElementos}{\nt{ElementoInicializador} \nt{RestoListaElementos}}\nobreak\par
{\leftskip=1.6em\relax\noindent\textup{Primeros} $=$ \{ \tl{!}, \tl{(}, \tl{+}, \tl{-}, \tl{charLiteral}, \tl{false}, \tl{floatLiteral}, \tl{idClase}, \tl{idMetVar}, \tl{intLiteral}, \tl{new}, \tl{null}, \tl{stringLiteral}, \tl{this}, \tl{true}, \tl{\{} \}\par}
\par\hangindent=2.2em\hangafter=1
\regla{RestoListaElementos}{\tl{,} \nt{ElementoInicializador} \nt{RestoListaElementos}}\nobreak\par
{\leftskip=1.6em\relax\noindent\textup{Primeros} $=$ \{ \tl{,} \}\par}
\par\hangindent=2.2em\hangafter=1
\regla{RestoListaElementos}{\eps}\nobreak\par
{\leftskip=1.6em\relax\noindent\textup{Primeros} $=$ \{ \eps \}\par}
\par\hangindent=2.2em\hangafter=1
\regla{ElementoInicializador}{\nt{Expresion}}\nobreak\par
{\leftskip=1.6em\relax\noindent\textup{Primeros} $=$ \{ \tl{!}, \tl{(}, \tl{+}, \tl{-}, \tl{charLiteral}, \tl{false}, \tl{floatLiteral}, \tl{idClase}, \tl{idMetVar}, \tl{intLiteral}, \tl{new}, \tl{null}, \tl{stringLiteral}, \tl{this}, \tl{true} \}\par}
\par\hangindent=2.2em\hangafter=1
\regla{ElementoInicializador}{\nt{InicializadorDeArreglo}}\nobreak\par
{\leftskip=1.6em\relax\noindent\textup{Primeros} $=$ \{ \tl{\{} \}\par}
\par\hangindent=2.2em\hangafter=1
\regla{DimensionesConTamanio}{\tl{[} \nt{Expresion} \tl{]} \nt{DimensionesConTamanioOpcional}}\nobreak\par
{\leftskip=1.6em\relax\noindent\textup{Primeros} $=$ \{ \tl{[} \}\par}
\par\hangindent=2.2em\hangafter=1
\regla{DimensionesConTamanioOpcional}{\nt{DimensionesConTamanio}}\nobreak\par
{\leftskip=1.6em\relax\noindent\textup{Primeros} $=$ \{ \tl{[} \}\par}
\par\hangindent=2.2em\hangafter=1
\regla{DimensionesConTamanioOpcional}{\eps}\nobreak\par
{\leftskip=1.6em\relax\noindent\textup{Primeros} $=$ \{ \eps \}\par}
\par\hangindent=2.2em\hangafter=1
\regla{ArgsActuales}{\tl{(} \nt{ListaExpsOpcional} \tl{)}}\nobreak\par
{\leftskip=1.6em\relax\noindent\textup{Primeros} $=$ \{ \tl{(} \}\par}
\par\hangindent=2.2em\hangafter=1
\regla{ListaExpsOpcional}{\nt{ListaExps}}\nobreak\par
{\leftskip=1.6em\relax\noindent\textup{Primeros} $=$ \{ \tl{!}, \tl{(}, \tl{+}, \tl{-}, \tl{charLiteral}, \tl{false}, \tl{floatLiteral}, \tl{idClase}, \tl{idMetVar}, \tl{intLiteral}, \tl{new}, \tl{null}, \tl{stringLiteral}, \tl{this}, \tl{true} \}\par}
\par\hangindent=2.2em\hangafter=1
\regla{ListaExpsOpcional}{\eps}\nobreak\par
{\leftskip=1.6em\relax\noindent\textup{Primeros} $=$ \{ \eps \}\par}
\par\hangindent=2.2em\hangafter=1
\regla{ListaExps}{\nt{Expresion} \nt{RestoListaExps}}\nobreak\par
{\leftskip=1.6em\relax\noindent\textup{Primeros} $=$ \{ \tl{!}, \tl{(}, \tl{+}, \tl{-}, \tl{charLiteral}, \tl{false}, \tl{floatLiteral}, \tl{idClase}, \tl{idMetVar}, \tl{intLiteral}, \tl{new}, \tl{null}, \tl{stringLiteral}, \tl{this}, \tl{true} \}\par}
\par\hangindent=2.2em\hangafter=1
\regla{RestoListaExps}{\tl{,} \nt{Expresion} \nt{RestoListaExps}}\nobreak\par
{\leftskip=1.6em\relax\noindent\textup{Primeros} $=$ \{ \tl{,} \}\par}
\par\hangindent=2.2em\hangafter=1
\regla{RestoListaExps}{\eps}\nobreak\par
{\leftskip=1.6em\relax\noindent\textup{Primeros} $=$ \{ \eps \}\par}
\endgroup

